% ============================================================
% Section 150: 广义 Virial 定理
% ============================================================
\section{量子力学：广义 Virial 定理}
\label{sec:hypervirial-theorem}

\begin{modelbox}[题目：广义 Virial 定理的证明]
在 Hermite 算符 $\hat A$ 的分立谱本征态下，证明 $\hat C=[\hat A,\hat B]$ 的期望值为零。
\end{modelbox}

\begin{tipbox}[考点快速分析]
\begin{itemize}
    \item \textbf{厄米算符}：$\int f^*(\hat A g)\,d\tau=\int(\hat A f)^*g\,d\tau$
    \item \textbf{本征方程}：$\hat A\varphi_n=a_n\varphi_n$，$a_n\in\mathbb{R}$
    \item \textbf{对易子期望}：$\langle[\hat A,\hat B]\rangle=\langle\hat A\hat B\rangle-\langle\hat B\hat A\rangle$
    \item \textbf{超 Virial 定理}：取 $\hat A=\hat H$ 时，$\langle[\hat H,\hat B]\rangle=0$
\end{itemize}
\end{tipbox}

\begin{derivationbox}[标准解题过程]
\textbf{已知条件}

$\hat A$ 为厄米算符，$\varphi_n$ 是其分立本征态，本征值 $a_n\in\mathbb{R}$：$\hat A\varphi_n=a_n\varphi_n$。$\hat B$ 为任意算符。

\textbf{证明过程}

计算对易子 $\hat C=[\hat A,\hat B]=\hat A\hat B-\hat B\hat A$ 在态 $\varphi_n$ 下的期望值：
\[
\langle\hat C\rangle_n=\int\varphi_n^*(\hat A\hat B-\hat B\hat A)\varphi_n\,d\tau.
\]

拆为两项。第一项利用 $\hat A$ 的厄米性：
\[
\int\varphi_n^*\hat A(\hat B\varphi_n)\,d\tau=\int(\hat A\varphi_n)^*(\hat B\varphi_n)\,d\tau=a_n\int\varphi_n^*\hat B\varphi_n\,d\tau.
\]

第二项直接利用本征方程：
\[
\int\varphi_n^*\hat B(\hat A\varphi_n)\,d\tau=a_n\int\varphi_n^*\hat B\varphi_n\,d\tau.
\]

因此
\begin{equation}
\boxed{\langle\hat C\rangle_n=a_n\langle\hat B\rangle_n-a_n\langle\hat B\rangle_n=0.}
\end{equation}
证毕。

\textbf{物理意义与讨论}

\textit{超 Virial 定理}：若取 $\hat A=\hat H$（哈密顿量），则结论变为：在能量本征态下，任意算符 $\hat B$ 与 $\hat H$ 的对易子期望值为零，即 $\langle[\hat H,\hat B]\rangle=0$。这在量子力学中被称为\textbf{超 Virial 定理}（Hypervirial Theorem）。

\textit{应用}：选择合适的 $\hat B$ 可推导出许多力学量平均值之间的关系，常用于计算势能平均值与动能平均值的关系（如经典的 Virial 定理 $\langle T\rangle=\frac{1}{2}\langle\mathbf{r}\cdot\nabla V\rangle$ 的量子对应）。
\end{derivationbox}

\begin{formulabox}[答案汇总]
\begin{center}
\begin{tabular}{ll}
\toprule
\textbf{结论} & \textbf{结果} \\
\midrule
广义 Virial 定理 & $\langle[\hat A,\hat B]\rangle_n=0$（在 $\hat A$ 的本征态下） \\
超 Virial 定理 & $\langle[\hat H,\hat B]\rangle=0$（在能量本征态下） \\
\bottomrule
\end{tabular}
\end{center}
\end{formulabox}

\begin{insightbox}[物理图像]
\begin{itemize}
    \item 证明的关键在于利用厄米算符的本征值为实数，以及厄米性带来的积分对称性。
    \item 超 Virial 定理是量子力学中极其有用的工具，通过选择不同的 $\hat B$（如 $\hat B=\mathbf{r}\cdot\mathbf{p}$、$\hat B=r^n$ 等），可以导出各种平均值关系式，避免直接求解波函数。
\end{itemize}
\end{insightbox}
